% Informe autocontenido. Compilar dos veces con pdfLaTeX (no requiere BibTeX).
\documentclass[11pt,a4paper]{article}
\usepackage[utf8]{inputenc}
\usepackage[T1]{fontenc}
\usepackage[spanish,es-nodecimaldot]{babel}
\usepackage[margin=2.3cm]{geometry}
\usepackage{lmodern,microtype}
\usepackage{amsmath,amssymb,bm}
\usepackage{booktabs,longtable,tabularx,array}
\usepackage{xcolor}
\usepackage{tikz}
\usetikzlibrary{arrows.meta,positioning}
\usepackage{enumitem}
\usepackage{listings}
\usepackage{xurl}
\usepackage[colorlinks=true,linkcolor=blue!50!black,citecolor=blue!50!black,urlcolor=blue!50!black]{hyperref}
\definecolor{ink}{RGB}{24,53,74}
\definecolor{pale}{RGB}{236,244,248}
\hypersetup{pdftitle={Arquitectura y fusión sensorial de AiGB UEx Navigation},pdfauthor={Informe de análisis del repositorio}}
\setlength{\parskip}{0.45em}
\setlength{\parindent}{0pt}
\setlength{\emergencystretch}{2em}
\setlist{nosep,leftmargin=*}
\newcommand{\code}[1]{\nolinkurl{#1}}
\newcommand{\file}[1]{\nolinkurl{#1}}
\newcommand{\R}{\mathbb{R}}
\newcommand{\trans}[2]{{}^{#1}\!T_{#2}}
\newcommand{\wrap}{\operatorname{wrap}_{[-\pi,\pi)}}
\newcommand{\clip}{\operatorname{clip}}
\lstset{basicstyle=\ttfamily\footnotesize,breaklines=true,columns=fullflexible,keepspaces=true,frame=single,backgroundcolor=\color{pale},showstringspaces=false}
\title{\textbf{Arquitectura y fusión sensorial de\newline AiGB\_UEx\_Navigation}\\[0.6em]\large Localización del Husky A300, percepción RGB-D y seguimiento humano con LiDAR y cámara ZED}
\author{Informe técnico basado en el código y la configuración del repositorio}
\date{14 de octubre de 2026}
\begin{document}
\maketitle

\begin{abstract}
Este informe reconstruye el funcionamiento de la versión de trabajo del repositorio
\code{AiGB_UEx_Navigation}. El sistema integra un robot Husky A300 simulado en Webots,
drivers RoboComp, transporte de datos mediante DDS y agentes que comparten un grafo
DSR. Hay tres mecanismos de fusión diferentes: un filtro de Kalman extendido que
localiza el robot mediante GPS, odometría de ruedas y orientación/velocidad angular
de la IMU; un seguimiento humano que combina esqueletos RGB-D con evidencia
geométrica del LiDAR; y una rejilla de ocupación que integra rayos LiDAR y profundidad
ZED para representar obstáculos. El código revisado no integra el LiDAR ni la
odometría visual de la ZED en el EKF de localización. La fusión humana tampoco es un
estimador conjunto de articulaciones y nube LiDAR: la cámara inicia las identidades
y estima la postura, mientras el LiDAR contribuye a la existencia y al desplazamiento
horizontal de personas previamente conocidas. Se documentan ecuaciones, marcos,
parámetros, fuentes concretas y límites de la implementación. Dos comprobaciones
aisladas existentes, del EKF y de la cinemática humana, se ejecutaron satisfactoriamente;
no se realizó una campaña de navegación o percepción en vivo.
\end{abstract}

\newpage
\setcounter{tocdepth}{1}
\tableofcontents
\newpage

\section{Alcance, método y trazabilidad}
\label{sec:scope}

El objeto de estudio es el \textbf{árbol de trabajo presente}, con los perfiles
\file{active_inference/utils/sub_husky.toml} y
\file{active_inference/utils/cognitive_husky.toml} como referencia de ejecución.
Se inspeccionaron los componentes, sus configuraciones, los productores y consumidores
de cada dato y las funciones que efectivamente se invocan. La identificación de un
algoritmo en un archivo o de una opción en TOML no se tomó por sí sola como prueba
de que esa ruta se ejecute.~\cite{launchsub,launchcog}

La referencia Git observada fue
\code{afab2d8cc17628b015118e930d40922f21cf6da5}. El árbol tenía modificaciones y
numerosos archivos eliminados localmente; por ello, ese identificador \emph{no}
describe por sí solo el contenido analizado. No se recuperaron componentes eliminados
ni se alteró la implementación. Las referencias finales dan rutas relativas,
funciones o intervalos de líneas del contenido revisado. Cuando un comentario y el
flujo ejecutable difieren, este informe prioriza el segundo.

Se distinguen cuatro niveles de evidencia:
\begin{itemize}
\item \textbf{Implementado}: comportamiento visible en las funciones y llamadas revisadas.
\item \textbf{Configurado}: valores del perfil Husky; no equivalen a una frecuencia o precisión medida.
\item \textbf{Comprobado aisladamente}: resultados de pruebas ejecutadas durante esta revisión.
\item \textbf{Pendiente de validación}: consecuencias plausibles del código, problemas potenciales y experimentos propuestos.
\end{itemize}

El README declara localización de aproximadamente 2 cm y navegación/seguimiento
funcionales en simulación. Se conserva esa afirmación como \textbf{declaración del
proyecto}, no como resultado experimental reproducido aquí. Los valores medidos
durante esta revisión se presentan separadamente en la sección~\ref{sec:tests}.
\cite{readme}

\section{Visión general del sistema}

\subsection{Organización del repositorio}
\begin{longtable}{p{0.27\linewidth}p{0.67\linewidth}}
\toprule
Directorio & Responsabilidad \\
\midrule
\endhead
\file{webots-agri/} & Modelos PROTO del Husky, terreno, palmeras y manzanos; mundos de simulación; generadores de hileras y malla de autofiltro; controlador de autocomprobación. \\
\file{webots-bridge/} & Accede a los dispositivos Webots, adapta coordenadas y mensajes, ofrece interfaces Ice y aplica órdenes de velocidad a las ruedas. También expone el GPS. \\
\file{robocomp-robolab/} & Drivers conservados: LiDAR, ZED, IMU y mando Xbox. Los tres primeros adaptan las fuentes al plano de datos sensorial. \\
\file{active_inference/robot_concept/} & Carga el modelo del robot en DSR, incorpora telemetría y anuncia los descriptores de los flujos sensoriales. \\
\file{active_inference/openfield_concept/} & EKF de localización planar, marco \code{field}, límite del área de trabajo y publicación de pose/covarianza. \\
\file{active_inference/retina/} & Percepción visual: segmentación de instancias y pose humana; publica máscaras y esqueletos. Contiene otras rutas deshabilitadas en el perfil Husky. \\
\file{active_inference/human_concept/} & Asociación temporal, modelo cinemático por persona, seguimiento con LiDAR, existencia y nodos \code{person_*}. \\
\file{active_inference/residual_concept/} & Ocupación de obstáculos a partir de LiDAR y profundidad; resta explicaciones geométricas al leer la rejilla. \\
\file{active_inference/controller/} & Planificación, seguimiento de ruta, tratamiento de obstáculos y mando de la base. \\
\file{active_inference/common/} & Transporte DDS, lectura temporal de transformaciones, existencia, geometría, huella, seguimiento y utilidades compartidas. \\
\file{active_inference/utils/} & Lanzadores, envío de objetivos y monitorización de procesos, comunicaciones y ventanas. \\
\bottomrule
\end{longtable}

La arquitectura procede de una adaptación de componentes para interiores. Quedan
nombres como \code{room}, \code{bpearl} o referencias a Shadow en comentarios y
funciones, pero el perfil agrícola utiliza \code{field}, un único LiDAR Helios y la
ZED frontal. \code{common/world_frame} prefiere el nodo de tipo \code{field} y conserva
\code{room} como alternativa. También existe \code{husky_apple.wbt}, además de los
mundos de palmeras descritos en el README; su presencia no demuestra que los
parámetros de seguimiento estén validados entre árboles.~\cite{readme,worldframe,worlds}

\subsection{Flujo completo}
\begin{figure}[htbp]
\centering
\begin{tikzpicture}[
 box/.style={draw=ink,rounded corners,fill=pale,align=center,font=\small,text width=4.1cm,minimum height=1.1cm},
 arr/.style={-{Latex[length=2mm]},thick,draw=ink}]
\node[box] (sim) at (0,0) {Webots: Husky y campo\\GPS, ruedas, IMU, LiDAR, RGB-D};
\node[box] (io) at (0,-1.9) {Bridge y drivers RoboComp\\Ice / IceStorm / DDS};
\node[box] (ekf) at (-4.7,-4) {\code{openfield_concept}\\GPS + ruedas + IMU $\to$ pose};
\node[box] (ret) at (0,-4) {\code{retina}\\RGB + profundidad $\to$ esqueleto};
\node[box] (res) at (4.7,-4) {\code{residual_concept}\\LiDAR + profundidad $\to$ rejilla};
\node[box] (hum) at (0,-6.1) {\code{human_concept}\\Esqueleto + LiDAR $\to$ personas};
\node[box] (ctl) at (0,-8.3) {\code{controller}\\Pose + rejilla + personas $\to$ mando};
\draw[arr] (sim)--(io);
\draw[arr] (io)--(ekf);
\draw[arr] (io)--(ret);
\draw[arr] (io)--(res);
\draw[arr] (ret)--(hum);
\draw[arr] (ekf)--node[left,font=\scriptsize]{marco común} (hum);
\draw[arr] (res.south) |- (ctl.east);
\draw[arr] (hum)--(ctl);
\draw[arr] (ekf.south) |- (ctl.west);
\draw[arr,dashed] (io.east) -- (2.5,-1.9) -- (2.5,-5.4) -- (hum.east)
 node[pos=.62,right,font=\scriptsize]{LiDAR};
\end{tikzpicture}
\caption{Dependencias funcionales principales. DSR contiene el estado compartido y las transformaciones; los datos sensoriales voluminosos circulan por el plano DDS. La pose también permite registrar la rejilla en el campo. El mando vuelve al bridge y a las ruedas.}
\label{fig:pipeline}
\end{figure}

\code{robot_concept} constituye la base común del grafo. El lanzador cognitivo
enumera después \code{openfield_concept}, \code{residual}, \code{controller},
\code{retina} y \code{human_concept}. El orden de la lista facilita el arranque,
pero no es sincronización de sensores: los agentes verifican disponibilidad del
grafo, de productores y de datos mediante sus ciclos y mecanismos de presencia.
\cite{launchcog,robotconfig}

\subsection{Qué significa «fusión» en este repositorio}
\begin{tabularx}{\linewidth}{p{0.20\linewidth}XXX}
\toprule
Problema & Entradas & Estado estimado & Mecanismo \\
\midrule
Localización & GPS, ruedas, IMU & $x,y,\theta$ del robot & EKF: predicción y correcciones. \\
Postura humana & RGB y profundidad ZED & Articulaciones y alineamiento del cuerpo & YOLO-pose, desproyección y ajuste cinemático. \\
Seguimiento humano & Esqueleto y LiDAR & Identidad, centro horizontal y existencia & Asociación geométrica, relevo de posición y evidencia en log-odds. \\
Obstáculos & LiDAR y profundidad ZED & Ocupación del campo & Integración de rayos en una rejilla común. \\
\bottomrule
\end{tabularx}

La pose del robot es una \textbf{entrada geométrica} de la percepción: permite situar
mediciones tomadas desde sensores móviles en el campo. No se encontró en el camino
activo una corrección recíproca del EKF basada en personas, máscaras o alineamiento
LiDAR. Por tanto, el sistema no debe describirse como SLAM LiDAR--visual--inercial ni
como un único filtro que fusiona todos los sensores.~\cite{ekf,openworker,humanlidar,resworker}

\section{Simulación, sensores y adquisición}

\subsection{Robot y escenarios}
El Husky tiene cuatro ruedas accionadas con cinemática \emph{skid-steer}; el giro
implica deslizamiento lateral. El modelo usa radio de rueda de 0.1651 m y semiancho
de vía de 0.2854 m en el bridge, con \code{Base.TrackScale = 1.875}. Este multiplicador
modifica la relación entre diferencia de velocidades de ruedas y giro; no elimina
la variación del deslizamiento con la maniobra. La elección explica que el EKF use
el giróscopo para predecir orientación.~\cite{proto,bridgeconfig,ekf}

El mundo de campo declara paso básico de 8 ms, referencia GPS
$(38.8794^\circ,-6.9707^\circ,186\,\mathrm{m})$ y norte en $+Y$. Dos peatones
recorren trayectorias predefinidas, a 0.5 y 0.6 m/s respectivamente. El archivo indica
que están destinados a percepción y carecen de objeto de colisión física. Por ello,
ver al robot atravesar o evitar gráficamente a un trabajador no constituye por sí
solo una prueba física de contacto o seguridad.~\cite{worlds}

\subsection{Especificación realmente utilizada}
\begin{longtable}{p{0.20\linewidth}p{0.73\linewidth}}
\toprule Sensor & Configuración observada \\
\midrule\endhead
LiDAR Helios & Montaje $(0,-0.25,0.95)$ m respecto a \code{body}. El PROTO genera 900 muestras horizontales y 50 capas sobre aproximadamente $110^\circ$ verticales, alcance 30 m y ruido nominal 0.02 m. El bridge conserva elevaciones entre $-55^\circ$ y $+15^\circ$: se retienen 32 capas, aproximadamente de $+14.6^\circ$ a $-55^\circ$. No deben confundirse las 50 capas generadas con las 32 publicadas tras el recorte. \\
ZED simulada & Montaje $(0,0.42,0.55)$ m, inclinación frontal descendente de aproximadamente $15^\circ$. Imagen $1280\times720$, campo horizontal de unos $110^\circ$. Profundidad mediante \code{RangeFinder}, límites 0.3--20 m. \\
GPS & Antena $(0,-0.25,1.02)$ m. Ruido configurado por eje de 0.02 m, correlación 0.9; actualización solicitada cada 100 ms. \\
IMU & Montaje $(0,0,0.2251)$ m en el grafo. El bridge lee giróscopo, acelerómetro y orientación de \code{InertialUnit}; el EKF utiliza $\omega_z$ y yaw. \\
\bottomrule
\end{longtable}
Fuentes: PROTO, modelo JSON del robot y configuración del bridge. Los valores son
propiedades del escenario revisado, no especificaciones garantizadas de equipos
reales.~\cite{proto,mounts,bridgeconfig}

Hay una discrepancia concreta: el descriptor configurado en el driver LiDAR conserva
\code{RangeMaxM = 100.0} y referencias textuales a Shadow, mientras el PROTO Husky
limita el sensor a 30 m. Esta diferencia afecta a la interpretación de metadatos;
el alcance efectivo simulado procede del dispositivo.~\cite{lidarconfig,proto}

\subsection{La profundidad simulada no se calcula por correspondencia estéreo}
\label{sec:simdepth}
La ruta Husky usa \code{Config.Simulated = true}. Webots produce una imagen de
\code{Camera} y una matriz de profundidad de \code{RangeFinder}; el bridge las
empaqueta conjuntamente, y \code{zed_camera} obtiene ese paquete mediante
\code{CameraRGBDSimple::getAll}. Esta ruta no reconstruye profundidad a partir de
dos imágenes ni ejecuta la estimación estéreo del SDK ZED. Es una \textbf{emulación
RGB-D de la interfaz ZED}.~\cite{bridgecamera,zeddriver,zedconfig}

El driver contiene otra rama para cámara real, con SDK ZED, profundidad
\code{NEURAL_PLUS} y seguimiento posicional. Sin embargo, la llamada que publicaría
la pose está comentada en \code{process_pose_data}. Además, la configuración activa es
simulada. La existencia de \code{enablePositionalTracking} en el driver no permite
afirmar que la odometría visual ZED localice el Husky en este perfil.~\cite{zeddriver}

\section{Comunicaciones, unidades, calibración y tiempo}

\subsection{Dos planos de información}
El plano DSR, configurado en dominio 0, representa robot, sensores, campo, personas,
rejilla y atributos. Los mensajes voluminosos se transmiten por DDS en dominio 7:
\code{rc/zed/rgb}, \code{rc/zed/depth}, \code{rc/lidar3d/points} y
\code{rc/imu/data}. Los drivers anuncian descriptores; \code{robot_concept} puede
retransmitirlos al grafo, y los consumidores descubren dominio, tópico y parámetros
del flujo a partir del nodo correspondiente. Los perfiles solicitan memoria
compartida y \emph{data sharing}. Esto reduce copias en el transporte, pero no implica
que no existan copias en el procesamiento de imágenes.~\cite{robotconfig,media,lidarconfig,zedconfig}

Ice se utiliza, entre otras cosas, para obtener RGB-D simulado, consultar el GPS y
mandar velocidades. IceStorm transporta publicaciones como la telemetría
\code{FullPoseEstimationPub}; \code{robot_concept} extrae de ella las velocidades de
ruedas y guarda aparte la pose de referencia del simulador.~\cite{bridgeconfig,robotworker,launchsub}

\subsection{Marcos de referencia y extrínsecas}
Se adopta la notación $\trans{A}{B}$ para la transformación que lleva coordenadas
del marco $B$ al marco $A$. El marco corporal usa $X$ hacia la derecha, $Y$ hacia
delante y $Z$ hacia arriba. El campo usa este, norte y arriba. La cámara en los
consumidores también usa $Y$ como eje de profundidad; no se debe aplicar sin
conversión una fórmula de cámara con $Z$ óptico.~\cite{ekfheader,mounts,graphpublisher}

El archivo \file{active_inference/robot_concept/husky.json} describe
\code{root -> husky -> body -> sensores}. El enlace \code{husky -> body} es
identidad. Las extrínsecas principales son:
\begin{align}
\trans{B}{L}&=\begin{bmatrix}I&(0,-0.25,0.95)^T\\0&1\end{bmatrix},\\
\trans{B}{C}&=\begin{bmatrix}R_x(-0.2618)&(0,0.42,0.55)^T\\0&1\end{bmatrix}.
\end{align}
El bridge y el driver convierten previamente los ejes nativos de Webots. Por eso,
copiar además la rotación de $90^\circ$ del dispositivo Webots en la arista DSR del
LiDAR duplicaría una conversión. El driver mantiene los puntos en el marco sensorial
normalizado y deja el montaje corporal en la cadena RT.~\cite{mounts,lidardriver,proto}

Para comparar geometrías tomadas en instantes distintos se busca:
\begin{equation}
 p^F_L=\trans{F}{B}(t_L)\trans{B}{L}p^L,
 \qquad p^F_C=\trans{F}{B}(t_C)\trans{B}{C}p^C.
\end{equation}
La comparación cámara--LiDAR se efectúa así en un marco común. En la ruta humana
no se encontró una proyección de todos los puntos LiDAR sobre las articulaciones 2D
ni una calibración extrínseca optimizada en línea: se usan los montajes declarados.
\cite{lidarreader,skeletonsource,humanlidar}

\subsection{Unidades y autofiltro}
El driver LiDAR convierte la representación interna en milímetros a metros antes
de publicar DDS. El autofiltro de malla se aplica antes de publicar; usa la geometría
del Husky y su montaje para retirar retornos del propio robot. Esa transformación
para el filtrado no significa que el flujo publicado quede en coordenadas corporales.
\cite{lidardriver,lidarconfig}

La profundidad simulada se transmite como \code{float32} en metros. Retina interpreta
\code{FORMAT_DEPTH_F32} directamente y convierte \code{Z16} de milímetros a metros.
Debe conservarse esta distinción al conectar hardware real: en la rama SDK revisada
se configura \code{sl::UNIT::MILLIMETER}, pero la publicación DDS copia los flotantes
como \code{FORMAT_DEPTH_F32} sin conversión visible. Es una \textbf{incompatibilidad
potencial de unidades en la ruta real}, no un fallo observado en la simulación.
\cite{zeddriver,zeddds,mediaplane,reszed}

\subsection{Sincronización efectiva y límites}
Retina empareja RGB y profundidad mediante \code{ZedFrameAligner}; el consumidor
solicita profundidad con una tolerancia de 20 ms respecto al RGB más reciente.
El procesamiento sigue una política de conservar el dato más nuevo: si el cálculo
no alcanza la frecuencia de adquisición, puede saltarse imágenes para limitar el
retraso.~\cite{aligner,mediaplane,zedsource}

El lector LiDAR común consulta transformaciones interpoladas al sello temporal del
barrido. El consumidor de esqueletos consulta \code{field <- zed} al instante
publicado en el nodo \code{skeleton}. No hay una barrera global que espere una pareja
LiDAR--cámara exacta: los agentes actualizan sus estados asincrónicamente.
\cite{lidarreader,skeletonsource}

Se emplean tres nociones temporales: reloj de simulación para integrar velocidades
simuladas, época en milisegundos para registrar sensores/RT, y reloj monótono para
caducidades de personas. El EKF usa \code{sim_stamp_ms} cuando existe. Sin esa
distinción, una simulación que avanza más despacio que el tiempo real produciría
desplazamientos incorrectos. La asociación de personas y la banda de movimiento,
en cambio, usan tiempo monótono del proceso; su sensibilidad al factor de tiempo
real merece una prueba específica.~\cite{openworker,humanworker,skeletonsource}

\section{Localización: GPS, ruedas e IMU mediante EKF}
\label{sec:ekf}

\subsection{Estado y entradas}
\code{OpenFieldEkf} mantiene
\begin{equation}
 \bm{x}=[x\;y\;\theta]^T,\qquad P\in\R^{3\times3}.
\end{equation}
Se trata de una pose planar: no estima altura, roll, pitch, velocidad como estado,
sesgo de giróscopo ni escalas de ruedas. La velocidad longitudinal $v_a$, lateral
$v_s$ y angular $\omega$ son entradas. La orientación $\theta$ gira el eje corporal
$+X$ respecto al este, de modo que el vector hacia delante es
$[-\sin\theta,\cos\theta]^T$.~\cite{ekf,ekfheader}

\code{read_odometry} lee velocidades publicadas por \code{robot_concept}.
\code{pump_imu} predice con $\omega_z$ del giróscopo; la velocidad de giro de ruedas
se conserva como lectura/diagnóstico, pero no se usa en la propagación del EKF.
El acelerómetro tampoco entra en este filtro.~\cite{openworker,robotworker}

\subsection{Predicción}
Con $\bar\theta=\theta_k+\omega_k\Delta t/2$, la implementación usa:
\begin{align}
 \Delta x&=(\cos\bar\theta\,v_s-\sin\bar\theta\,v_a)\Delta t,\\
 \Delta y&=(\sin\bar\theta\,v_s+\cos\bar\theta\,v_a)\Delta t,\\
 \bm{x}_{k+1}^{-}&=[x_k+\Delta x,\ y_k+\Delta y,\ \wrap(\theta_k+\omega_k\Delta t)]^T.
\end{align}
El Jacobiano implementado y la propagación de incertidumbre son:
\begin{equation}
 F=\begin{bmatrix}1&0&-\Delta y\\0&1&\Delta x\\0&0&1\end{bmatrix},
 \qquad P^- = FPF^T+GQG^T,
\end{equation}
\begin{equation}
 G=\begin{bmatrix}
 -\sin\bar\theta\Delta t&\cos\bar\theta\Delta t&0\\
 \cos\bar\theta\Delta t&\sin\bar\theta\Delta t&0\\
 0&0&\Delta t
 \end{bmatrix},\quad
 Q=\operatorname{diag}(\sigma_a^2,\sigma_s^2,\sigma_\omega^2)/\Delta t.
\end{equation}
Las desviaciones son densidades dependientes del movimiento:
\begin{align}
 \sigma_a&=\sqrt{0.01^2+(0.05|v_a|)^2},\\
 \sigma_s&=\sqrt{0.005^2+(0.10|v_a\omega|)^2},\\
 \sigma_\omega&=\sqrt{0.002^2+(0.01|\omega|)^2}.
\end{align}
Así, $GQG^T$ crece proporcionalmente a $\Delta t$, y el ruido lateral aumenta en
maniobras que combinan avance y giro. Son coeficientes configurados; no se halló en
esta revisión una calibración experimental nueva que los identificase.~\cite{ekf,openconfig}

\subsection{GPS y brazo de palanca}
El GPS observa la antena, no el origen del robot. Para
$\bm{\ell}=(0,-0.25)^T$ m, el modelo de observación es:
\begin{equation}
 h_G(\bm{x})=\begin{bmatrix}x\\y\end{bmatrix}+R(\theta)\bm{\ell},
 \qquad H_G=\left[I_2\quad \frac{\partial R(\theta)\bm{\ell}}{\partial\theta}\right].
\end{equation}
El brazo de palanca acopla orientación y posición: una rotación desplaza la antena
aunque el origen corporal permanezca fijo. La inicialización resta el brazo rotado
al primer GPS disponible y propaga la incertidumbre del yaw a la posición y a sus
covarianzas cruzadas.~\cite{ekf}

Con $R_G=0.05^2I_2$, se calcula innovación, ganancia y corrección:
\begin{align}
 \bm{r}&=\bm{z}_G-h_G(\bm{x}^-),&S&=H_GP^-H_G^T+R_G,\\
 K&=P^-H_G^TS^{-1},&\bm{x}^+&=\bm{x}^-+K\bm{r}.
\end{align}
La covarianza se actualiza en forma de Joseph:
\begin{equation}
 P^+=(I-KH_G)P^-(I-KH_G)^T+KR_GK^T.
\end{equation}
La distancia de Mahalanobis se devuelve como diagnóstico; no se utiliza como
rechazo estadístico de GPS atípicos. El valor configurado de 0.05 m es superior al
ruido nominal del dispositivo simulado para compensar parcialmente la correlación
de sus muestras. No se estima un estado explícito de ruido GPS correlacionado.
\cite{ekf,openconfig}

El bridge convierte WGS84 a UTM y resta la referencia del mundo para entregar
coordenadas locales. Aunque el proyecto las denomina ENU, la operación concreta
implementada es una diferencia de coordenadas UTM; no debe confundirse con una
conversión geodésica general mediante ECEF y plano tangente local.~\cite{bridgegps}

\subsection{Corrección de orientación}
La IMU proporciona una medición absoluta de yaw:
\begin{equation}
 r_\theta=\wrap(z_\theta-\theta^-),\quad
 H_\theta=[0\;0\;1],\quad R_\theta=0.005^2.
\end{equation}
La actualización es escalar y también usa forma de Joseph. El perfil limita estas
correcciones a una cada 100 ms de reloj sensorial; la predicción continúa con cada
IMU nueva. El objetivo es evitar tratar todas las orientaciones de alta frecuencia
como observaciones independientes.~\cite{ekf,openworker,openconfig}

En simulación, el yaw procede de \code{InertialUnit}; no es una implementación propia
de fusión acelerómetro--giróscopo--magnetómetro. En un robot real será necesario
definir qué estimador y qué referencia absoluta suministran ese yaw.~\cite{bridgeimu}

\subsection{Publicación y dependencias}
Cada ciclo del agente lee odometría, procesa IMU, consulta GPS y publica. El período
configurado es 20 ms; el GPS se consulta cada 50 ms, aunque su dispositivo se
actualiza cada 100 ms. La deduplicación del GPS compara igualdad exacta de las
coordenadas, ya que la interfaz consultada no aporta sello temporal del fix. Esta
estrategia evita reutilizar un valor repetido, pero no distingue fixes diferentes
con iguales coordenadas.~\cite{openworker,openconfig}

La arista almacenada tiene padre robot y destino \code{field}; contiene la inversa
de la pose estimada:
\begin{equation}
 \trans{B}{F}=\begin{bmatrix}R^T&-R^T\bm{t}\\0&1\end{bmatrix}.
\end{equation}
El código transforma la covarianza con el Jacobiano de inversión y la inserta en
los índices $x,y,\mathrm{yaw}$ de una matriz $6\times6$. Publica también velocidad
y su incertidumbre. El resto de agentes pide la transformación en el sentido que
necesita mediante la API RT. Los ceros en las otras dimensiones expresan aquí una
representación planar; no son una medición de precisión perfecta en altura o
inclinación.~\cite{opengraph}

La referencia del supervisor se usa para visualizar y registrar el error, buscando
la estimación correspondiente al timestamp de esa referencia. No entra en las
correcciones del EKF. Desactivar GPS en el visor deja navegación por integración de
ruedas/giróscopo con crecimiento de incertidumbre; desactivar yaw elimina la
corrección absoluta de orientación.~\cite{openworker,robotworker}

\section{Percepción visual: de RGB-D a personas}

\subsection{Segmentación y pose son ramas diferentes}
Retina ejecuta sobre RGB una etapa de segmentación configurada con
\code{yolo26x-seg.onnx} y una etapa de pose con \code{yolo26l-pose.onnx}. Son nombres
de modelos solicitados por el perfil; no se certifican aquí sus pesos ni su
rendimiento. La segmentación acepta \code{person} y otras clases, con confianza
mínima 0.25 e IoU de supresión 0.45. La pose humana usa confianza mínima de persona
0.30, confianza mínima de articulación 0.30 y decimación 1.~\cite{retinaconfig,retinaworker}

Las máscaras se publican en \code{masks}; los esqueletos en \code{skeleton}. La ruta
de \code{human_concept} consume este último. No se encontró que la máscara de
segmentación \code{person} sea requisito de aceptación del esqueleto o que se use
para seleccionar sus muestras de profundidad. Por tanto, la fusión humana debe
describirse como YOLO-pose más profundidad, no como una validación obligatoria entre
dos detectores.~\cite{graphpublisher,skeletonsource}

\subsection{Reconstrucción 3D de las articulaciones}
Para cada articulación 2D $(u,v)$ aceptada se toman profundidades válidas de una
ventana de semiancho 2 píxeles, es decir, $5\times5$. El código usa el elemento de
orden central del conjunto de profundidades, una mediana robusta a valores extremos.
Con $d$ esa profundidad, la desproyección es:
\begin{equation}
 \bm{p}^{C}=\left[
 \frac{(u-c_x)d}{f_x},\ d,\ \frac{(c_y-v)d}{f_y}
 \right]^T,\qquad c_x=W/2,\quad c_y=H/2.
\end{equation}
Las focales proceden del flujo RGB-D. Un keypoint sin profundidad positiva y finita,
fuera de imagen o con confianza insuficiente se publica como NaN y confianza cero.
La ventana no está condicionada por una máscara humana: en un brazo delgado o un
borde puede incluir fondo, incluso si la articulación 2D es correcta.~\cite{graphpublisher}

Los 17 puntos COCO se reordenan al formato BODY\_18. Se añade el cuello como punto
medio tridimensional de los hombros cuando ambos son válidos; su confianza es la
mínima de las dos. El nodo \code{skeleton} incluye contador de publicación, timestamp,
número de cuerpos, IDs, $18\times3$ coordenadas por cuerpo y 18 confianzas. Las
confianzas se transportan en $[0,1]$ y se convierten a $[0,100]$ en el consumidor.
\cite{graphpublisher,skeletonsource}

\subsection{Registro en el campo y asociación temporal}
\code{DsrSkeletonSource::poll} transforma los esqueletos de cámara a campo y
\code{assign_track_ids} asigna identidades persistentes según proximidad espacial.
El centro se obtiene primero de las dos caderas; si no están disponibles, del
cuello; finalmente, de la media de articulaciones válidas. Para una pista con edad
$\Delta t$, la distancia máxima de asociación es:
\begin{equation}
 g_{\mathrm{cam}}=0.35+2.0\Delta t\quad [\mathrm{m}].
\end{equation}
Se ordenan los pares admisibles por distancia y se asignan vorazmente sin reutilizar
un cuerpo ni una pista. Las pistas sin actualización durante 3 s se retiran. Un
cuerpo sin asociación a menos de 0.5 m de otro ya asignado puede eliminarse como
duplicado.~\cite{skeletonsource}

Este mecanismo es asociación geométrica, sin descriptor de apariencia ni
reidentificación aprendida. La exclusión uno a uno evita compartir un track en un
frame, pero no garantiza conservar la identidad al cruzarse trabajadores, durante
oclusiones largas o cuando sus centros se confunden.~\cite{skeletonsource}

\section{Modelo cinemático humano e inferencia}

\subsection{Qué estima el modelo}
Cada \code{person_*} mantiene un modelo BODY\_18 con 11 grados de libertad: tres
ángulos por hombro, una flexión por codo y tres parámetros de cuerpo inferior
($lb_x$, $lb_z$, $lb_{roll}$). No es un modelo de articulación completa de las piernas.
El modelo canónico usa $Y$ vertical y se alinea al mundo mediante Kabsch; este marco
interno no debe confundirse con el marco corporal del robot, donde $Z$ es vertical.
\cite{kinematic}

Para cada persona, las longitudes de segmentos se calibran con las observaciones
válidas: la primera medición inicializa y las siguientes se incorporan mediante
media exponencial con peso 0.05. La estimación de postura requiere al menos tres
articulaciones utilizables para observar; el núcleo necesita además al menos tres
anclajes válidos para el alineamiento. Los anclajes configurados son cuello,
hombros y caderas. \code{MinValid = 12} se carga en configuración y sus comentarios
lo relacionan con acciones de observación, pero no se encontró un uso ejecutable de
ese campo en el seguimiento revisado; no es el mínimo efectivo del ajuste.
\cite{humanfitter,humanconfig,vfe}

\subsection{Ajuste por precisión y aproximación de Laplace}
El estimador mantiene una media $\bm{\mu}\in\R^{11}$ y una precisión aproximada
$\Lambda$. Una expresión explicativa de su objetivo es:
\begin{equation}
 \mathcal{F}(\bm{\theta})=
 \frac{1}{2M}\sum_{j\in\mathcal{V}}\pi_j
 \|R(\bm{\theta})f_j(\bm{\theta})+\bm{t}(\bm{\theta})-\bm{y}_j\|^2
 +\frac{1}{2}\|\bm{\theta}-\bm{\mu}_{prev}\|^2_{\Lambda_{dyn}}
 +\mathcal{P}(\bm{\theta}).
\end{equation}
$f_j$ es la cinemática directa, $R,\bm{t}$ el alineamiento Kabsch, $\mathcal{V}$ el
conjunto válido, y $\mathcal{P}$ agrupa restricciones de simetría, límites, cruce de
brazos y postura neutra. La expresión resume los términos del código; no añade
observaciones LiDAR al ajuste articular.~\cite{vfe}

Con confianza $c_j\in[0,1]$ se configura:
\begin{equation}
 \sigma_j=0.02+(0.15-0.02)(1-c_j),\qquad \pi_j=\sigma_j^{-2},
 \qquad \Lambda_{dyn}=0.25^{-2}I.
\end{equation}
La interpretación probabilística de esta precisión es una modelización elegida:
la confianza de la red no es, por sí misma, una varianza métrica calibrada. El prior
dinámico se aplica por actualización con esa matriz fija; no debe atribuirse a este
núcleo el mismo ruido por densidad temporal del EKF.~\cite{humanconfig,vfe}

Se realizan dos pasos Gauss--Newton. El Jacobiano se obtiene mediante diferencias
finitas centrales, de paso $10^{-3}$, a través de cinemática y Kabsch. Para el
residuo ya ponderado:
\begin{equation}
 H\simeq \frac{J^TJ}{M}+\Lambda_{dyn}+\lambda I+H_{neutral},\qquad
 \bm{\mu}\leftarrow\bm{\mu}-H^{-1}\bm{g},\qquad \Lambda\leftarrow H,
 \quad\lambda=10^{-3}.
\end{equation}
Algunos priors contribuyen solo al gradiente; no se usa el Hessiano exacto de todos
los términos. El núcleo aplica además límites articulares y rechazo de innovaciones
excesivas. Los términos internos de penalización de velocidad y aceleración están
desactivados en el perfil revisado; un controlador de salida limita el movimiento
del esqueleto publicado.~\cite{vfe,humanconfig,humanfitter}

\subsection{Salida e incertidumbre}
El esqueleto publicado usa límites de velocidad/aceleración y suavizado de la
traslación y de la rotación mediante cuaterniones, con peso 0.2 a la pose nueva.
La posición del nodo persona se toma de la pelvis. El seguimiento LiDAR modifica
esta geometría ya ajustada; no estima ángulos de brazos a partir de la nube.
\cite{humanfitter,humangraph,humanlidar}

La covarianza RT de la persona es una aproximación heurística:
\begin{equation}
 \sigma^2_{\mathrm{pos}}=\clip(0.02\,\operatorname{tr}(\Lambda^{-1}),10^{-4},1).
\end{equation}
Se asigna ese valor a las tres diagonales de posición y 1 a las de orientación.
No es la propagación completa de profundidad, calibración, pose del robot y
alineamiento global. Además, el valor llamado \emph{free energy} en parte de la
telemetría del agente se alimenta con \code{mean_l2}, un error medio en metros;
no debe publicarse como si fuera el valor de $\mathcal{F}$ de la ecuación anterior.
\cite{humangraph,humanfitter}

\section{Fusión LiDAR--ZED para seguimiento humano}
\label{sec:humanfusion}

\subsection{Nacimiento y reparto de responsabilidades}
La cámara inicia las hipótesis: \code{scaffold_missing_person_nodes} crea una persona
desde un esqueleto con pelvis válida. El LiDAR recorre instancias ya existentes;
no contiene en esta ruta un clasificador que cree personas nuevas a partir de
clusters. Una persona que nunca ha visto la cámara puede producir ocupación
geométrica, pero no obtiene por este mecanismo una identidad humana.~\cite{humangraph,humanlidar}

El LiDAR aporta dos funciones:
\begin{enumerate}
\item \textbf{Existencia}: retornos compatibles confirman; rayos que atraviesan el
volumen esperado aportan evidencia contraria; oclusión o ausencia de rayos no aporta
evidencia geométrica.
\item \textbf{Posición horizontal}: cuando no hay visión reciente, los puntos
asociados desplazan el esqueleto conocido y la pista de identidad.
\end{enumerate}
No hay una media ponderada simultánea de posiciones ZED y LiDAR con covarianza
conjunta. La implementación es una \textbf{fusión de evidencia y seguimiento con
prioridad visual y relevo geométrico}.~\cite{humanlidar,humanworker}

\subsection{Selección de puntos y puerta espacial}
Se seleccionan puntos en una banda global $z\in[1.20,1.85]$ m. Esta banda pretende
observar torso/cabeza por encima de las palmeras jóvenes. No sigue la anatomía de
cada esqueleto ni estima automáticamente el suelo bajo cada trabajador.
\cite{humanlidar,humanconfig}

Sea $r=0.25$ m el radio corporal, $\rho$ la distancia al sensor y $\Delta t$ el tiempo
desde la última posición actualizada. El código construye:
\begin{align}
 \sigma_{fix}&=\begin{cases}0.002\rho^2,&\text{última posición de cámara},\\0.05,&\text{última posición de LiDAR},\end{cases}\\
 \sigma_{turn}&=|\omega_{robot}|\,0.05\,\rho,\\
 \sigma&=\min(\sigma_{fix}+\sigma_{turn}+2.0\Delta t,\,1.0),\\
 g_L&=\min(r+2\sigma,\,1.0).
\end{align}
Estas magnitudes son una envolvente heurística en metros; la suma lineal no es una
propagación completa de covarianzas independientes. El término cuadrático con
distancia modela incertidumbre de profundidad de cámara; el término de giro
representa registro imperfecto durante latencia.~\cite{humanlidar,humanconfig}

Cada punto de la banda se asigna a la persona más cercana cuya puerta lo admite.
Esto evita que el mismo punto confirme simultáneamente dos centroides durante esta
asignación, aunque no garantiza una identidad correcta si las puertas se solapan.
Se admite incluso un solo punto para actualizar el centro, una decisión adaptada
a la escasez de capas sobre una persona lejana.~\cite{humanlidar}

\subsection{Actualización del centro}
Si $\bm{m}$ es la media de puntos asociados y $\bm{o}$ el origen LiDAR en el plano,
el centro propuesto es:
\begin{equation}
 \bm{c}_L=\bm{m}+\frac{r}{2}\frac{\bm{m}-\bm{o}}{\|\bm{m}-\bm{o}\|}.
\end{equation}
Se añade $r/2=0.125$ m hacia el interior porque los retornos representan la superficie
visible. Esta corrección es geométrica y fija; un único punto lateral o una postura
inusual puede sesgarla.~\cite{humanlidar}

Si la cámara se consideró observada desde el barrido previo, se conserva su centro
y se registra la separación \code{lidar_cam_gap_m}. En caso contrario, se traslada
en $XY$ el esqueleto de salida, el último ajuste y su traslación suavizada. Se marca
el fix como LiDAR y se llama a \code{refresh_track} para que la próxima detección
visual pueda asociarse a la misma identidad. No se actualiza una velocidad humana
latente ni una distribución cinemática de movimiento independiente.
\cite{humanlidar,humanworker,skeletonsource}

\subsection{Fusión probabilística de existencia}
La existencia se representa con log-odds:
\begin{equation}
 L=\log\frac{P(E)}{1-P(E)},\qquad P(E)=\frac{1}{1+e^{-L}}.
\end{equation}
El nacimiento usa $L_0=2$ y se limita $L$ al intervalo $[-4,4]$. Los parámetros del
modelo de detección son $p_D=0.85$ y $p_C=0.05$. Una observación de cámara aceptada
añade:
\begin{equation}
 a=\log\frac{p_D}{p_C}=\log17\simeq2.833.
\end{equation}
La cámara no incorpora en esta ruta una actualización negativa explícita por cada
frame sin esqueleto; esa ausencia se trata a través del LiDAR y del tiempo sin
evidencia.~\cite{humanlidar,humanworker,existence}

Para el LiDAR se intersectan rayos con una caja de torso de planta $0.5\times0.5$ m
y la banda vertical indicada. El código separa evidencia de retorno superficial
$e_{occ}$, retorno interior $e_{int}$ y espacio atravesado $e_{free}$. En este modelo
humano, tanto superficie como interior cuentan a favor. Un rayo detenido antes del
volumen no demuestra ausencia.~\cite{humanlidar,existence}

La evidencia negativa se atenúa si el centro es incierto:
\begin{align}
 b&=\log\frac{1-p_D}{1-p_C}\simeq-1.846,\\
 w_{free}&=\left(\frac{r}{r+\sigma_{carve}}\right)^2,\\
 \Delta L_L&=a\tanh\left(
 \frac{(e_{occ}+e_{int})a+w_{free}e_{free}b}{a}\right).
\end{align}
La saturación evita que muchos puntos correlacionados de un barrido cuenten como
muchas confirmaciones independientes. Cuando no hay retornos asociados,
$\sigma_{carve}$ es la envolvente anterior; con centro LiDAR se usa 0.05 m más el
término de giro. En la llamada a \code{carve_box}, la incertidumbre de superficie
pasada es 0.05 m; la envolvente ampliada interviene después en $w_{free}$, no como
un ensanchamiento idéntico de toda la geometría del ray tracing.~\cite{humanlidar}

En esta ruta no se configura la correlación entre frames de
\code{ExistenceBelief}; su valor por defecto es cero. Por tanto, la saturación dentro
del barrido sí está presente, pero no debe afirmarse que la aplicación humana
descorrelaciona automáticamente toda evidencia temporal ni toda dependencia entre
cámara y LiDAR. La creencia combina las aportaciones secuencialmente y aplica el
límite tras cada integración.~\cite{existence,humanfitter,humanlidar}

\subsection{Persistencia y eliminación}
Una persona se elimina si $P(E)<0.2$, o si pasan más de 2 s sin evidencia de cámara
ni barridos que alcancen su volumen. Este segundo mecanismo es un timeout explícito.
Así, la oclusión mantiene el log-odds al no haber medición, pero \textbf{no garantiza
persistencia indefinida}: el timeout puede eliminar la identidad. «No observación»
y «espacio libre» no son equivalentes, aunque ambos puedan terminar retirando un
nodo por mecanismos distintos.~\cite{humanworker,humanconfig}

El orden real de un ciclo es: leer esqueletos, crear nodos, preparar las observaciones,
ejecutar el paso LiDAR sobre instancias existentes, ajustar/publicar cada persona y
eliminar las ausentes. Por ello, la prueba de visión reciente del paso LiDAR se basa
en marcas de ajustes anteriores; un esqueleto recién leído se incorpora al ajuste
después de ese paso. Esta es otra razón para no describir el sistema como una
actualización multimodal perfectamente simultánea.~\cite{humanworker}

\section{Fusión geométrica de obstáculos y rejilla residual}
\label{sec:grid}

\subsection{Ruta activa}
El ciclo de \code{residual_concept} requiere un nuevo barrido LiDAR. Inicializa una
rejilla de resolución 0.05 m a partir de los límites del campo, con margen, e integra
el LiDAR y la nube obtenida de la profundidad ZED. Los rayos de cámara usan su propio
origen, necesario para no atribuir a la cámara el espacio libre visible desde otro
sensor.~\cite{resworker,grid,resconfig}

La ZED se restringe a profundidades de 0.3--5 m para esta ruta; el límite no afecta
del mismo modo al esqueleto humano, cuya desproyección no comparte ese recorte de
\code{ZedBoost}. Se submuestrea y voxeliza la nube de cámara, y se filtra
infraestructura con una tolerancia que depende de distancia:
\begin{equation}
 \sigma_{infra}(d)=0.03+0.006d^2.
\end{equation}
El ajuste del plano del suelo y el modelo de incertidumbre reducen la probabilidad
de convertir suelo ruidoso en obstáculos. Con el Husky no existe Bpearl; la función
de selección del suelo contiene un retorno alternativo al barrido disponible
cuando no obtiene puntos Bpearl.~\cite{resworker,resconfig,zedboost}

\subsection{Acumulación de evidencia}
La rejilla acumula hits y evidencia de espacio libre por celda; al cerrar un ciclo
se aplica una única actualización agregada, con precedencia del hit cuando hay
conflicto. Tiene histéresis, memoria de estabilidad, representación de soporte en
altura, tratamiento de la huella propia y mecanismos de olvido. El esquema
conceptual $L_{i,k+1}=L_{i,k}+\Delta L_{i,k}$ no describe por sí solo todos esos
factores; en particular, los pesos cambian con movimiento, alcance y certeza de pose.
\cite{grid,resworker}

Cuando está activada \code{Grid.PosePrecision}, un peso usado por la rejilla es de
la forma:
\begin{equation}
 w_{pose}(\rho)=\frac{s_{ref}^2}{s_{ref}^2+\sigma_p^2+\rho^2\sigma_\theta^2}.
\end{equation}
Esto vincula explícitamente localización y mapeado: una pose incierta reduce la
autoridad de los rayos para confirmar y borrar celdas. No convierte esos rayos en
correcciones de pose del robot.~\cite{grid,resworker}

El producto activo es un nodo de rejilla con celdas y datos derivados para el
controlador. Se excluyen explicaciones de infraestructura y objetos al leer los
residuales, sin borrar por ello toda la evidencia sensorial subyacente. La semántica
densa de suelo solo actúa si existe su producto; la etapa ADE20K está deshabilitada
en la configuración Husky, aunque permanezcan opciones subordinadas de publicación
con valor verdadero.~\cite{resworker,retinaworker,retinaconfig}

\subsection{Código conservado que no debe confundirse con la ruta activa}
Existen clases de DBSCAN, ajuste de cajas residuales y la opción
\code{ZedBoost.FeedDetection = true}. Sin embargo, en el \code{compute()} revisado
la salida se obtiene de la rejilla; no se llama a la cadena antigua de creación y
ajuste de instancias residuales. La variable \code{zed_detection_enabled} se carga
pero no aparece utilizada en ese camino. Los comentarios de
\code{residual_zed_boost.h} y de \code{observe()} que describen puntos añadidos a
clusters o cajas son insuficientes para sostener que ese es el algoritmo activo.
\cite{resworker,resconfigloader,zedboost}

\section{Cómo repercute la percepción en la navegación}
El controlador toma la pose del campo, el polígono de trabajo, la rejilla residual,
observaciones LiDAR y objetos del grafo. Su planificador de rejilla usa A* con
comprobación de la huella del robot; el perfil tiene celda de 0.06 m y ruta continua.
El modo de control seleccionado es \code{plain}, con \code{BandEnabled = true};
los comentarios heredados sobre MPPI no describen por sí solos el controlador
configurado actualmente. Los límites declarados son 0.7 m/s de avance y 1.0 rad/s
de giro.~\cite{controllerconfig,planner,trajectory}

Cada nodo \code{person} añade un obstáculo poligonal de 12 lados que aproxima un
disco de radio 0.8 m alrededor de su posición. El planificador sigue aplicando además
la huella del robot. Esto permite evitar una persona incluso si su geometría no
está bien representada en la banda de ocupación residual. El radio es un parámetro
de evitación, no una distancia de seguridad certificada ni una predicción de la
trayectoria humana.~\cite{obstacletracker,controllerconfig}

Debe distinguirse \textbf{seguimiento perceptivo} de personas de \textbf{seguirlas
con la base}. La ruta principal demostrada por el código revisado mantiene su
identidad, publica postura/posición y las incorpora como obstáculos. El agente
humano también propone acciones de observación valoradas mediante la precisión
posterior del grado de libertad peor observado. Eso no demuestra la existencia de una política completa
de acompañamiento de trabajadores ni de reconocimiento de su tarea agrícola.
\cite{humanworker,epistemic,readme}

\section{Parámetros esenciales para reproducir el perfil}
\begin{longtable}{p{0.38\linewidth}p{0.22\linewidth}p{0.32\linewidth}}
\toprule Parámetro & Valor & Interpretación \\
\midrule\endhead
Bridge: \code{Period.Compute} & 8 ms & Paso solicitado de adquisición/actuación. \\
Bridge: \code{Period.Sensors} & 50 ms & Muestreo solicitado de sensores pesados. \\
Bridge: \code{GPS.Period} & 100 ms & Muestreo del receptor simulado. \\
Bridge: \code{FullPose.PublishPeriod} & 100 ms & Publicación de telemetría de ruedas. \\
Openfield: \code{Period.Compute} & 20 ms & Drena IMU, consulta GPS y publica pose. \\
Openfield: \code{GpsSigmaM} & 0.05 m & Desviación de observación del EKF. \\
Openfield: \code{YawSigmaRad} & 0.005 rad & Desviación de yaw absoluto. \\
Retina: \code{Period.Compute} & 25 ms & Ciclo del agente; el worker tiene su propio coste. \\
Retina: \code{HumanPose.hold_ms} & 500 ms & Caché de poses ante fallos transitorios. \\
Retina: \code{depth_patch} & 2 px & Semiancho para muestrear profundidad. \\
Human: \code{Period.Compute} & 50 ms & Ciclo de seguimiento. \\
Human: \code{BandZMin/Max} & 1.20 / 1.85 m & Banda LiDAR global. \\
Human: \code{FollowGateMaxM} & 1.0 m & Límite de puerta espacial. \\
Human: \code{DeathTimeoutS} & 2.0 s & Tiempo máximo sin evidencia. \\
Residual: \code{Period.Compute} & 100 ms & Ciclo con requisito de nuevo LiDAR. \\
Residual: \code{ZedBoost.MaxDepthM} & 5.0 m & Recorte de profundidad para ocupación. \\
Controller: \code{ControlMode} & \code{plain} & Seguimiento seleccionado. \\
Controller: \code{PersonRadiusM} & 0.8 m & Radio del obstáculo humano. \\
\bottomrule
\end{longtable}
Fuentes: perfiles activos. Las frecuencias inversas de estos períodos son nominales.
El paso Webots, la carga gráfica, la planificación de hilos y el coste ONNX pueden
reducir la frecuencia efectiva. El driver ZED configurado a 30 FPS no crea 30 frames
nuevos por segundo si el productor simulado entrega menos.~\cite{bridgeconfig,openconfig,retinaconfig,humanconfig,resconfig,controllerconfig,zedconfig}

\section{Limitaciones y hallazgos que deben figurar en el reporte}
\label{sec:limits}

\subsection{No existe fusión LiDAR--ZED en el EKF del Husky}
La localización activa depende del GPS y del yaw absoluto para corregir deriva.
No se observó scan matching, corrección por landmarks visuales, odometría estéreo
ni cierre de bucle conectados a \code{OpenFieldEkf}. La disponibilidad de LiDAR y
cámara no equivale a usarlos para estimar la pose.~\cite{ekf,openworker}

\subsection{Modelo planar en un terreno tridimensional}
El terreno puede tener camellones e irregularidades, pero la pose global publicada
es SE(2), con altura y roll/pitch nulos en la arista dinámica. Si el robot se inclina,
los montajes estáticos y el yaw no bastan para registrar perfectamente la nube y
la profundidad en 3D. El efecto sobre la banda del torso y el suelo es una hipótesis
geométrica que requiere medición, no un fallo cuantificado en esta revisión.
\cite{worlds,opengraph}

\subsection{Reutilización de poses antiguas como si fueran nuevas}
\code{YoloHumanProcessor::detect_poses} conserva \code{last_poses_} hasta 500 ms
después de la última detección si el detector no devuelve personas.
\code{PoseStage::run} marca \code{poses_fresh = run_now}, que significa «se ejecutó
el detector», no «se obtuvo una detección nueva». El agente publica después esos
keypoints usando la profundidad y el timestamp del bundle actual.~\cite{yolohuman,posestage,retinaworker}

Por tanto, en una pérdida transitoria de detección, una pose 2D retenida puede
volver a desproyectarse con profundidad nueva y contarse como evidencia visual
reciente. Esta posibilidad contradice la lectura literal de «retención solo para
el visor» de algunos comentarios. No se midió su frecuencia en ejecución. Una
corrección futura debería separar datos retenidos para visualización de observaciones
con adquisición real nueva.~\cite{graphpublisher,humanworker}

\subsection{Transformación ausente en el consumidor humano}
Si \code{DsrSkeletonSource::poll} no obtiene \code{field <- zed}, utiliza identidad
y devuelve puntos de cámara; escribe un aviso, pero no descarta necesariamente el
cuerpo. Si el campo existe pero falla la cadena temporal RT, esos puntos podrían
acabar interpretados como coordenadas globales por el creador de personas. Es un
riesgo de registro visible en el código; no se reproduce aquí un caso concreto.
\cite{skeletonsource,humangraph,humanworker}

\subsection{Asincronía y reutilización de profundidad en ocupación}
\code{integrate_zed_into_grid} llama a \code{pump()} pero comprueba
\code{has_depth()}, sin exigir que la llamada haya devuelto un frame nuevo. Además,
consulta la transformación de cámara con \code{current_ts_}, el timestamp del
LiDAR, en lugar de emplear explícitamente el de profundidad. Una profundidad retenida
puede reintegrarse al llegar otros barridos; con movimiento o retraso entre sensores
puede introducir evidencia repetida o registro desplazado. No debe afirmarse que
esta ruta posee la misma sincronización RGB-D de 20 ms de retina.
\cite{resworker,reszed}

\subsection{Supuestos del seguimiento LiDAR}
La banda 1.20--1.85 m y la aceptación de un solo punto explotan un escenario con
vegetación baja. Ramas, árboles, vehículos u objetos altos pueden producir retornos
compatibles. Un trabajador agachado puede salir de la banda, y un terreno elevado
altera su altura global. La clasificación semántica inicial procede de la cámara;
las confirmaciones LiDAR no son un clasificador humano generalizable sin validación.
La puerta máxima de 1 m y el timeout de 2 s también limitan las pérdidas tolerables.
\cite{humanlidar,humanconfig,worlds}

\subsection{Incertidumbres parciales y dependencia de la simulación}
La covarianza de persona es un proxy; la puerta LiDAR es una envolvente heurística;
el ajuste humano convierte confianza de red a ruido métrico por una regla elegida.
No existe un posterior conjunto que incluya estado del robot, extrínsecas, IDs,
articulaciones y todas sus correlaciones. Tampoco la emulación RGB-D reproduce por
sí misma fallos reales de correspondencia estéreo. Estas limitaciones impiden
trasladar directamente resultados simulados a una ZED y un LiDAR físicos.
\cite{humangraph,humanlidar,vfe,zeddriver}

\subsection{GPS retrasado y odometría retenida}
La interfaz GPS utilizada no entrega timestamp del fix; el filtro aplica la
corrección al estado tras drenar IMU. No se observó una actualización retardada con
rebobinado del estado. Asimismo, \code{read_odometry} conserva la última velocidad
si no hay atributo nuevo, y \code{predict} sigue usándola mientras llega IMU. El
timeout primario del agente vigila IMU; el tratamiento de pérdidas de ruedas y
retrasos GPS debe ensayarse separadamente.~\cite{openworker,openconfig}

\subsection{Documentación y opciones heredadas}
El README del núcleo humano todavía presenta como pendiente la integración con
DSR, aunque esa integración sí existe. La configuración LiDAR cita un alcance
heredado; \code{FeedDetection} no demuestra una ruta activa de clusters; el perfil
del controlador mantiene comentarios MPPI con modo \code{plain}. Por ello, las
afirmaciones de este informe se apoyan en llamadas ejecutables además de documentación.
\cite{humancorereadme,lidarconfig,resworker,controllerconfig}

\section{Validación realizada y evaluación propuesta}
\label{sec:tests}

\subsection{Comprobaciones ejecutadas durante la revisión}
Se compiló y ejecutó el \code{self_test()} existente de \code{OpenFieldEkf}, aislado
del resto del sistema. La prueba genera 20\,000 pasos a 8 ms, equivalentes a 160 s
de movimiento sintético, con tramos rectos, pivotes, arcos y paradas. Tras los
primeros 5 s, comprueba error máximo de posición inferior a cinco desviaciones GPS,
error máximo de yaw inferior a $1^\circ$ y positividad de la covarianza final.
\cite{ekf}

\begin{center}
\begin{tabular}{lr}
\toprule Resultado del test sintético & Valor impreso \\
\midrule
Estado & PASS \\
Máximo error de posición & 0.035 m \\
Máximo error de yaw & $0.317^\circ$ \\
Covarianza final definida positiva & Sí \\
Desviación final horizontal agregada & 0.011 m \\
Desviación final de yaw & 0.0017 rad \\
\bottomrule
\end{tabular}
\end{center}

Estos son resultados redondeados de esa secuencia artificial con semilla fija,
\textbf{no} una precisión medida en el campo Webots. El test no reproduce toda la
asincronía de procesos, la correlación GPS configurada en el PROTO o el terreno.
Incluso sus comentarios de tasas deben interpretarse con cuidado: el bucle aplica
GPS y yaw cada 12 pasos en fases distintas; a 8 ms son intervalos de 96 ms para cada
uno.~\cite{ekf}

También se compiló el núcleo humano independiente de DSR/Qt y se ejecutó su
\code{sanity_checks} mediante CTest: \textbf{1/1 prueba pasó}. Verifica propiedades
básicas de cinemática, Kabsch y ajuste de una postura sintética sin ruido;
no certifica detección YOLO, seguimiento de IDs ni
fusión LiDAR. No se ejecutó la prueba de paridad con el antiguo prototipo Python.
\cite{humantest,humancmake}

No se lanzó el stack sensoriomotor ni se mandaron movimientos al robot. Tampoco se
generaron medidas nuevas de precisión/recall, error por articulación, tasa de
cambio de identidad, latencia global o colisiones. El informe es principalmente
una auditoría del flujo de implementación, complementada por las dos pruebas
aisladas anteriores.

\subsection{Protocolo recomendado para un reporte experimental}
Las siguientes pruebas son propuestas, no resultados:
\begin{enumerate}
\item \textbf{Localización}: repetir trayectorias rectas, giros sobre sí mismo y
arcos en los mundos plano y de campo. Comparar EKF completo, GPS desactivado y yaw
desactivado mediante los interruptores existentes. Registrar RMSE, percentiles y
máximos de posición/yaw contra la referencia temporalmente alineada.
\item \textbf{Consistencia}: analizar innovaciones GPS y, si se instrumentan,
NIS/NEES. Un error pequeño no demuestra que $P$ esté bien calibrada. Inyectar
retardos y pérdidas independientes de GPS, ruedas e IMU.
\item \textbf{Ablación humana}: comparar \code{HumanLidar.Enabled=true/false} con
la misma trayectoria. Medir duración de pistas, cambios de identidad, falsos
nacimientos, pérdidas al salir del campo visual y error del centro frente a
trayectorias de referencia.
\item \textbf{Escenarios difíciles}: cruce de dos trabajadores, persona agachada,
oclusiones, giro del robot y rangos crecientes. Separar «persona nunca detectada»
de «persona conocida mantenida por LiDAR».
\item \textbf{Timestamps}: registrar edades RGB, profundidad, LiDAR, pose y
publicación de persona; comprobar la retención de keypoints y profundidad descrita
en la sección~\ref{sec:limits}.
\item \textbf{Ocupación}: comparar rejilla LiDAR sola con LiDAR más ZED,
desactivando \code{ZedBoost.FeedGrid} para la primera condición. Medir falsos
obstáculos de suelo, cobertura, tiempo de borrado y efecto sobre rutas.
\item \textbf{Transferencia}: antes de ensayar hardware, verificar unidades,
extrínsecas y fuente del yaw absoluto; después evaluar con escenas y oclusiones
reales. Los peatones sin colisión física no bastan como evaluación de seguridad.
\end{enumerate}

Para reproducibilidad deben guardarse el mundo, configuraciones, pesos ONNX y sus
hashes, versión del árbol de trabajo, semilla cuando exista, factor de tiempo real,
hardware de inferencia y timestamps. Los CSV existentes de localización, ajuste
humano y ocupación aportan parte de esta evidencia; los IDs y la pose de referencia
humana necesitarían comprobarse o instrumentarse para una evaluación completa.
\cite{openworker,humanworker,resworker}

\section{Construcción y ejecución del proyecto}
El proyecto requiere dependencias externas: RoboComp, Cortex/DSR con el tipo
\code{field}, Webots, Eigen, Qt, FastDDS, ONNX Runtime y los modelos solicitados por
retina. El README especifica compilación por componente y limita el paralelismo a
ocho trabajos. No hay una compilación monolítica del repositorio que sustituya a
todas esas instalaciones.~\cite{readme}

Los pasos documentados para el perfil simulado son:
\begin{lstlisting}
# En cada componente que deba construirse:
cmake -B build
make -C build -j8

# Con IceStorm disponible y el mundo Husky abierto en Webots:
cd active_inference/utils
python3 subcognitive.py sub_husky.toml

# En otra terminal, desde el mismo directorio:
python3 cognitive.py cognitive_husky.toml

# Objetivo de ejemplo en coordenadas del campo; requiere
# habilitar el control mediante Run en el perfil actual:
python3 send_goal.py X Y
\end{lstlisting}
Los lanzadores contienen rutas basadas en
\file{~/robocomp/AiGB_UEx_Navigation}. Los modelos ONNX no se versionan según el
README. La ausencia de un modelo, del descriptor DDS o de una transformación puede
desactivar etapas o impedir resultados aun cuando el proceso se haya iniciado.
\cite{launchsub,launchcog,readme,posestage}

El presente documento es autocontenido y no requiere archivo bibliográfico externo:
\begin{lstlisting}
cd docs
pdflatex -interaction=nonstopmode -halt-on-error informe_fusion_sensorial.tex
pdflatex -interaction=nonstopmode -halt-on-error informe_fusion_sensorial.tex
\end{lstlisting}

\section{Redacción sintética utilizable en un technical report}
El sistema desarrollado organiza la percepción y navegación mediante agentes
RoboComp que comparten un grafo DSR. La localización del Husky A300 se obtiene con
un EKF planar cuya predicción combina velocidad de ruedas y velocidad angular de
IMU, y cuyas correcciones utilizan GPS y yaw absoluto. La observación GPS incorpora
el desplazamiento de la antena respecto al origen corporal. La pose y su covarianza
permiten expresar datos sensoriales en el marco global del campo.

La detección humana parte de keypoints 2D obtenidos por YOLO-pose en la imagen ZED;
su profundidad local permite reconstruir un esqueleto BODY\_18 tridimensional.
Después del registro en el campo y la asociación temporal, un modelo cinemático
de 11 grados de libertad ajusta la postura mediante pasos Gauss--Newton y una
aproximación de Laplace. El LiDAR complementa a la visión sobre personas ya
detectadas: confirma o refuta su existencia mediante rayos en una banda de torso,
y actualiza su centro horizontal cuando la cámara pierde observación reciente.
Las aportaciones de existencia se combinan en log-odds y la posición se actualiza
mediante prioridad visual y relevo LiDAR, sin un filtro articular multimodal único.

En paralelo, una rejilla residual integra LiDAR y profundidad RGB-D para representar
obstáculos. El controlador combina esa ocupación con la pose del robot y obstáculos
derivados de las personas. En el perfil simulado, la profundidad se obtiene de un
RangeFinder de Webots; no se ha demostrado aquí la equivalencia de rendimiento con
una ZED física. La arquitectura implementa por tanto localización GPS--inercial--
odométrica y percepción geométrica/semántica complementaria, con limitaciones de
sincronización, asociación e incertidumbre que deben evaluarse experimentalmente.

% REFERENCIAS_GENERADAS: rutas y localizadores del arbol de trabajo revisado.
\clearpage
Las fuentes siguientes son código y configuraciones locales consultados el 14 de octubre de 2026; no son referencias bibliográficas externas. Los localizadores se refieren al árbol de trabajo analizado, y las huellas muestran los primeros 16 dígitos hexadecimales de SHA-256. Un archivo conservado no implica que todas sus rutas estén activas.

\begin{thebibliography}{99}
\bibitem{readme} Descripción general, dependencias y lanzamiento.
\file{README.md}.
Archivo completo (86 líneas). Huella SHA-256 abreviada: \code{e53ac2a50fbf1131}.

\bibitem{launchsub} Perfil de lanzamiento sensoriomotor.
\file{active_inference/utils/sub_husky.toml}.
Archivo completo (46 líneas). Huella SHA-256 abreviada: \code{131d62e3696e99a4}.

\bibitem{launchcog} Perfil de lanzamiento cognitivo.
\file{active_inference/utils/cognitive_husky.toml}.
Archivo completo (43 líneas). Huella SHA-256 abreviada: \code{4c793a6f79160d90}.

\bibitem{worldframe} Selección del marco global.
\file{active_inference/common/world_frame/world_frame.h}.
Archivo completo (88 líneas). Huella SHA-256 abreviada: \code{91a84e75fbc4842c}.

\bibitem{worlds} Mundo agrícola y peatones; véanse también husky\_test.wbt y husky\_apple.wbt.
\file{webots-agri/worlds/husky_field.wbt}.
Localizadores de inicio: líneas 15, 56. Huella SHA-256 abreviada: \code{160086b558d88a63}.

\bibitem{proto} Geometría y sensores simulados.
\file{webots-agri/protos/HuskyA300.proto}.
Localizadores de inicio: líneas 39, 135, 167. Huella SHA-256 abreviada: \code{eb9de552ab1e018f}.

\bibitem{bridgeconfig} Perfil Husky del bridge.
\file{webots-bridge/etc/config_husky}.
Archivo completo (203 líneas). Huella SHA-256 abreviada: \code{d96ec91601d80de3}.

\bibitem{mounts} Modelo DSR del robot y extrínsecas.
\file{active_inference/robot_concept/husky.json}.
Archivo completo (494 líneas). Huella SHA-256 abreviada: \code{8843802c663bbd2e}.

\bibitem{robotconfig} Perfil del productor del grafo.
\file{active_inference/robot_concept/etc/config_husky.toml}.
Archivo completo (240 líneas). Huella SHA-256 abreviada: \code{57b1261263ec2843}.

\bibitem{robotworker} Telemetría de ruedas y referencia del supervisor.
\file{active_inference/robot_concept/src/specificworker.cpp}.
Localizadores de inicio: líneas 2126. Huella SHA-256 abreviada: \code{c67af76ee10bd696}.

\bibitem{bridgecamera} Adquisición de RGB y profundidad.
\file{webots-bridge/src/specificworker.cpp}.
Localizadores de inicio: líneas 361. Huella SHA-256 abreviada: \code{6940a92fced0e74a}.

\bibitem{bridgegps} Conversión y consulta GPS.
\file{webots-bridge/src/specificworker.cpp}.
Localizadores de inicio: líneas 2123, 2180, 2225. Huella SHA-256 abreviada: \code{6940a92fced0e74a}.

\bibitem{bridgeimu} Lectura de IMU del simulador.
\file{webots-bridge/src/specificworker.cpp}.
Localizadores de inicio: líneas 1963. Huella SHA-256 abreviada: \code{6940a92fced0e74a}.

\bibitem{lidarconfig} Montaje, filtro y descriptor LiDAR.
\file{robocomp-robolab/components/hardware/laser/lidar3d_dds/etc/config_helios_husky_webots.toml}.
Archivo completo (277 líneas). Huella SHA-256 abreviada: \code{1fc2163bc7e3b70b}.

\bibitem{lidardriver} Conversión de nube, autofiltro y publicación.
\file{robocomp-robolab/components/hardware/laser/lidar3d_dds/src/specificworker.cpp}.
Localizadores de inicio: líneas 419, 601, 695. Huella SHA-256 abreviada: \code{ecd4c5ef8eb85867}.

\bibitem{zedconfig} Perfil ZED simulado.
\file{robocomp-robolab/components/hardware/camera/zed_camera/etc/configHuskyWebot.toml}.
Archivo completo (58 líneas). Huella SHA-256 abreviada: \code{44333298e7a24157}.

\bibitem{zeddriver} Ramas SDK/simulación, profundidad y publicación de pose.
\file{robocomp-robolab/components/hardware/camera/zed_camera/src/specificworker.cpp}.
Localizadores de inicio: líneas 205, 342, 483, 600, 660. Huella SHA-256 abreviada: \code{7d70f57475feb156}.

\bibitem{zeddds} Formato de profundidad DDS.
\file{robocomp-robolab/components/hardware/camera/zed_camera/src/dds_publisher.cpp}.
Localizadores de inicio: líneas 124. Huella SHA-256 abreviada: \code{9761d4b518916920}.

\bibitem{media} Descriptores, tópicos y fábricas de suscriptores.
\file{active_inference/common/media_transport/media_transport.h}.
Localizadores de inicio: líneas 466, 503. Huella SHA-256 abreviada: \code{d17b5e7d00bfa0df}.

\bibitem{lidarreader} Lectura de barridos y transformación al instante de captura.
\file{active_inference/common/media_transport/lidar_plane_reader.cpp}.
Localizadores de inicio: líneas 134, 161, 191. Huella SHA-256 abreviada: \code{7b084b17b8b4bbed}.

\bibitem{mediaplane} Unidades de profundidad y emparejamiento RGB-D.
\file{active_inference/retina/src/media_plane_source.cpp}.
Localizadores de inicio: líneas 163, 442, 458. Huella SHA-256 abreviada: \code{5312aebccd09b6b6}.

\bibitem{aligner} Buffer temporal de RGB y profundidad.
\file{active_inference/retina/src/zed_frame_aligner.cpp}.
Localizadores de inicio: líneas 34. Huella SHA-256 abreviada: \code{816702efa93783fc}.

\bibitem{zedsource} Selección del frame más nuevo y pose de cámara.
\file{active_inference/retina/src/zed_source.cpp}.
Localizadores de inicio: líneas 26. Huella SHA-256 abreviada: \code{98a3361b9490d4e2}.

\bibitem{ekf} Inicialización, ecuaciones EKF y prueba sintética.
\file{active_inference/openfield_concept/src/openfield_ekf.cpp}.
Localizadores de inicio: líneas 32, 48, 85, 109, 129. Huella SHA-256 abreviada: \code{e28bce92da1d9a80}.

\bibitem{ekfheader} Estado, marcos y parámetros del EKF.
\file{active_inference/openfield_concept/src/openfield_ekf.h}.
Localizadores de inicio: líneas 31, 51. Huella SHA-256 abreviada: \code{71c2c5aa9f6d3b57}.

\bibitem{openworker} Orquestación de sensores y evaluación contra referencia.
\file{active_inference/openfield_concept/src/specificworker.cpp}.
Localizadores de inicio: líneas 203, 223, 278, 297, 380. Huella SHA-256 abreviada: \code{2c82d557713c7fbd}.

\bibitem{openconfig} Parámetros de localización y campo.
\file{active_inference/openfield_concept/etc/config.toml}.
Archivo completo (86 líneas). Huella SHA-256 abreviada: \code{5f09dd9efa600ab3}.

\bibitem{opengraph} Creación de field, inversión de pose y publicación RT.
\file{active_inference/openfield_concept/src/openfield_scene_graph.cpp}.
Localizadores de inicio: líneas 24, 70. Huella SHA-256 abreviada: \code{f425aceb1fb14f0c}.

\bibitem{retinaconfig} Modelos y etapas visuales del perfil Husky.
\file{active_inference/retina/etc/config_husky.toml}.
Archivo completo (541 líneas). Huella SHA-256 abreviada: \code{f5c1c80dbe469ff1}.

\bibitem{retinaworker} Construcción de etapas y publicación de resultados.
\file{active_inference/retina/src/specificworker.cpp}.
Localizadores de inicio: líneas 193, 209, 1716. Huella SHA-256 abreviada: \code{3579c370351fe4f1}.

\bibitem{graphpublisher} Publicación y desproyección BODY 18.
\file{active_inference/retina/src/graph_publisher.cpp}.
Localizadores de inicio: líneas 911, 925. Huella SHA-256 abreviada: \code{85e5be6c00e75118}.

\bibitem{yolohuman} Detector y retención de poses.
\file{active_inference/retina/src/yolo_human.cpp}.
Localizadores de inicio: líneas 283. Huella SHA-256 abreviada: \code{ba649483bab297e7}.

\bibitem{posestage} Etapa de pose y semántica de poses\_fresh.
\file{active_inference/retina/src/pose_stage.cpp}.
Localizadores de inicio: líneas 24. Huella SHA-256 abreviada: \code{2d470d065959c32b}.

\bibitem{skeletonsource} Registro, asociación y refresco de IDs.
\file{active_inference/human_concept/src/skeleton_source.cpp}.
Localizadores de inicio: líneas 108, 225, 244. Huella SHA-256 abreviada: \code{b7d6ae1731ed2d61}.

\bibitem{humanconfig} Parámetros humanos y LiDAR.
\file{active_inference/human_concept/etc/config.toml}.
Archivo completo (106 líneas). Huella SHA-256 abreviada: \code{c86def28a08049d4}.

\bibitem{humanfitter} Instancias, observación, inferencia y suavizado.
\file{active_inference/human_concept/src/human_fitter.cpp}.
Localizadores de inicio: líneas 75, 107, 136. Huella SHA-256 abreviada: \code{aaa498cc3cd09936}.

\bibitem{kinematic} Modelo BODY 18 y definición de 11 grados de libertad.
\file{active_inference/human_concept/cpp/core/human_kinematic_model.h}.
Localizadores de inicio: líneas 17, 109. Huella SHA-256 abreviada: \code{fadbc7d0da7adb09}.

\bibitem{vfe} Kabsch, precisiones, Gauss--Newton y aproximación de Laplace.
\file{active_inference/human_concept/cpp/core/vfe_inference.cpp}.
Localizadores de inicio: líneas 23, 208. Huella SHA-256 abreviada: \code{3690da77f58fe1d7}.

\bibitem{humangraph} Nacimiento de personas, esqueleto y covarianza proxy.
\file{active_inference/human_concept/src/human_scene_graph.cpp}.
Localizadores de inicio: líneas 52, 122. Huella SHA-256 abreviada: \code{234a11a833e18dde}.

\bibitem{humanlidar} Fusión de existencia y seguimiento geométrico.
\file{active_inference/human_concept/src/human_lidar_presence.cpp}.
Localizadores de inicio: líneas 52, 62. Huella SHA-256 abreviada: \code{bd7dfda4caa528e6}.

\bibitem{humanworker} Orden de actualización, eliminación y acciones de observación.
\file{active_inference/human_concept/src/specificworker.cpp}.
Localizadores de inicio: líneas 250, 299, 327, 358, 416. Huella SHA-256 abreviada: \code{be3bf9a780f5d5a8}.

\bibitem{existence} Ray casting, saturación y creencia de existencia.
\file{active_inference/common/existence_belief/existence_belief.h}.
Localizadores de inicio: líneas 103, 112, 392, 499. Huella SHA-256 abreviada: \code{63d39fb3dee8e749}.

\bibitem{epistemic} Propuesta de vista y ganancia esperada de información.
\file{active_inference/human_concept/src/epistemic_planner.cpp}.
Localizadores de inicio: líneas 16. Huella SHA-256 abreviada: \code{fa80c3a6471a327f}.

\bibitem{resworker} Ruta activa de ocupación, integración ZED y código residual conservado.
\file{active_inference/residual_concept/src/specificworker.cpp}.
Localizadores de inicio: líneas 672, 1115, 1378, 1828. Huella SHA-256 abreviada: \code{f9f5a039acd83f22}.

\bibitem{resconfig} Perfil de rejilla y profundidad.
\file{active_inference/residual_concept/etc/config_husky.toml}.
Archivo completo (519 líneas). Huella SHA-256 abreviada: \code{abdec61d37361854}.

\bibitem{resconfigloader} Carga de FeedDetection y otras opciones.
\file{active_inference/residual_concept/src/residual_config.cpp}.
Localizadores de inicio: líneas 104. Huella SHA-256 abreviada: \code{a61e0c8a6a6c2cdb}.

\bibitem{reszed} Retención y unidades del frame de profundidad.
\file{active_inference/residual_concept/src/residual_zed_ingestor.cpp}.
Localizadores de inicio: líneas 69. Huella SHA-256 abreviada: \code{af5655569a0a1b01}.

\bibitem{zedboost} Geometría de desproyección y límites de nube ZED.
\file{active_inference/residual_concept/src/residual_zed_boost.h}.
Localizadores de inicio: líneas 37, 63. Huella SHA-256 abreviada: \code{eb3d46a0b418088d}.

\bibitem{grid} Integración de rayos, precisión de pose y consolidación de celdas; parámetros en el .h homónimo.
\file{active_inference/residual_concept/src/residual_occupancy_grid.cpp}.
Localizadores de inicio: líneas 486, 263, 524. Huella SHA-256 abreviada: \code{3338460e7f8fbfd8}.

\bibitem{controllerconfig} Modo plain, ruta, límites y obstáculos humanos.
\file{active_inference/controller/etc/config_husky.toml}.
Localizadores de inicio: líneas 122, 289, 533, 555. Huella SHA-256 abreviada: \code{e7084e2b78bf44d8}.

\bibitem{planner} Planificación de rejilla y comprobación de huella.
\file{active_inference/controller/src/grid_planner.cpp}.
Localizadores de inicio: líneas 653. Huella SHA-256 abreviada: \code{fef90402ef02dfdb}.

\bibitem{trajectory} Seguimiento y modificación local de trayectoria.
\file{active_inference/controller/src/trajectory_controller.cpp}.
Archivo completo (1347 líneas). Huella SHA-256 abreviada: \code{37ea2982f25cbca5}.

\bibitem{obstacletracker} Obstáculos derivados de nodos person.
\file{active_inference/controller/src/controller_obstacle_tracker.cpp}.
Localizadores de inicio: líneas 1180. Huella SHA-256 abreviada: \code{9e58ff7aedab5e69}.

\bibitem{humancorereadme} Documentación del núcleo y tareas históricas pendientes.
\file{active_inference/human_concept/cpp/README.md}.
Archivo completo (57 líneas). Huella SHA-256 abreviada: \code{13697f290e9b654d}.

\bibitem{humancmake} Construcción aislada y registro de sanity\_checks.
\file{active_inference/human_concept/cpp/CMakeLists.txt}.
Localizadores de inicio: líneas 32. Huella SHA-256 abreviada: \code{ccf38a57d46989b6}.

\bibitem{humantest} Comprobaciones de cinemática, Kabsch e inferencia sintética.
\file{active_inference/human_concept/cpp/test/sanity_checks.cpp}.
Localizadores de inicio: líneas 28. Huella SHA-256 abreviada: \code{90382a4d1ed33025}.

\end{thebibliography}
\end{document}
